Las tres prácticas de laboratorio permitieron construir, de forma incremental, un entendimiento comparado de cómo se materializa la arquitectura clásica de RPC —cliente, stub del cliente, \emph{runtime} de RPC, stub del servidor y servidor— en tres tecnologías con niveles de abstracción distintos. La Práctica 1 expuso, sin abstracciones adicionales, cada componente de la arquitectura clásica mediante CGI sobre nginx, donde el marshalling de parámetros es manual y la semántica de invocación depende enteramente del diseño del script. La Práctica 2 mostró la misma arquitectura resuelta de forma nativa por el lenguaje Java a través de RMI, donde el stub se genera dinámicamente y la semántica \emph{at-most-once} se ofrece por defecto en el propio \emph{runtime}. La Práctica 3 integró ambos modelos en una arquitectura heterogénea orquestada por nginx como \emph{gateway}, reproduciendo a escala reducida el patrón habitual de los sistemas distribuidos contemporáneos, en los que coexisten distintos protocolos de invocación remota dentro de una misma solución.
